% !TEX root = ../monograph.tex
\chapter{Polymer expansions and connected correlations}
\label{ch08:chapter}

An effective density is useful only if it remains local after integration.
Its logarithm should therefore be expressed as a sum of terms supported on
connected sets, with a bound that does not grow with the surrounding volume.
The cluster expansion provides this passage from local activities to connected
terms. Its convergence is a standard result of statistical mechanics. The
work specific to the gauge theory is to put the activities produced by the
block integration within its hypotheses.

This chapter states the abstract results in the form needed for that
application. It does not prove another cluster-expansion theorem. We use
Koteck\'y--Preiss convergence\cite{KoteckyPreiss1986}, in the explicit rooted
form of Ueltschi\cite[Theorem~1]{Ueltschi2004}, and explain how support weights,
activity derivatives and large-field constraints enter. The distinction
between a general cluster estimate and a bound on the actual gauge-theory
activities will remain explicit throughout.

\section{Activities, supports and incompatibility}
\label{ch08:activities}

Let $\Lambda$ be a finite set of lattice sites. A \emph{polymer label}
$\gamma$ belongs to a finite or countable set $\mathcal P$, has a nonempty
support $X_\gamma\subset\Lambda$, and carries a complex number
$z_\gamma$, called its \emph{activity}. Different labels may have the same
support. The labels can retain information about an integration sector,
derivative or local Taylor coefficient; they are not identified solely by
their supports.

A symmetric, reflexive relation $\gamma\not\sim\eta$ specifies
\emph{incompatibility}. Reflexivity means $\gamma\not\sim\gamma$: two occurrences of the same
label are incompatible. We assume
\begin{equation}\label{ch08:overlap}
 \gamma\not\sim\eta\quad\Longrightarrow\quad
 X_\gamma\cap X_\eta\ne\varnothing.
\end{equation}
When incompatibility is defined by overlap of supports, the implication
is an equivalence. We use this particular rule for the constrained gas
in Section~\ref{ch08:forced}.

The polymer partition function is
\begin{equation}\label{ch08:gas}
 Z_{\rm pol}(z)=\sum_{\substack{A\subset\mathcal P\text{ finite}\\
                         A\text{ pairwise compatible}}}
                       \prod_{\gamma\in A}z_\gamma.
\end{equation}
The empty collection contributes one. This sum is a polynomial for finitely
many labels. For countably many labels it is well defined whenever
$\sum_\gamma|z_\gamma|<\infty$, since its absolute sum is at most
$\exp(\sum_\gamma|z_\gamma|)$. The cluster criterion below implies this
summability in finite volume. Complex or signed activities are allowed;
they need not be probability weights.

We keep track of spatial size with a nonnegative function $c(X)$ on
nonempty subsets of $\Lambda$. The needed properties are
\begin{equation}\label{ch08:cover-properties}
 |X|\le h\,c(X),\qquad
 c(X\cup Y)\le c(X)+c(Y)\quad\text{if }X\cap Y\ne\varnothing,
\end{equation}
where $h>0$ is fixed. One choice is the number of sites, with $h=1$.
A useful multiscale choice is the minimum number of radius-$B$ lattice
balls whose union is connected and contains $X$. In that case
$h=(2B+1)^4$ is an upper bound on the size of a ball, including on a torus.
The symbol $c$ will always denote this support-size function in this chapter,
not a field coefficient.

For $\alpha\ge0$ define the site norm
\begin{equation}\label{ch08:site-norm}
 \|z\|_\alpha=
 \sup_{x\in\Lambda}\sum_{\gamma:x\in X_\gamma}
                  |z_\gamma|e^{\alpha c(X_\gamma)}.
\end{equation}
It counts every polymer at each site in its support. A norm which counts
only the chosen anchor of a polymer is different. Passing from an anchored
bound to \eqref{ch08:site-norm} requires a support-size estimate or a loss
in the exponential weight. This distinction matters when checking a
cluster criterion for an effective density.

\section{The logarithm and its convergence criterion}
\label{ch08:criterion}

For a tuple $\boldsymbol\gamma=(\gamma_1,\ldots,\gamma_n)$,
form the graph on the \emph{occurrences} $1,\ldots,n$, joining distinct
vertices whenever their labels are incompatible. Define
\begin{equation}\label{ch08:ursell}
 \varphi^T(\boldsymbol\gamma)=
 \sum_{\substack{H\text{ connected spanning graph}\\
                   H\text{ uses incompatibility edges}}}
                 (-1)^{|E(H)|}.
\end{equation}
For $n=1$ its value is one; for a disconnected incompatibility graph it
is zero. Repeated labels are allowed in this tuple, even though the gas
\eqref{ch08:gas} contains each label at most once. They are necessary in
its logarithm: for a gas with one label, $Z_{\rm pol}=1+z$ and
$\log Z_{\rm pol}=z-z^2/2+\cdots$. The $n!$ below compensates for
ordering the occurrences; it is not part of the definition
\eqref{ch08:ursell}.

\begin{theorem}[Standard polymer convergence criterion]
\label{ch08:kp}
Suppose $A_\gamma\ge0$ and
\begin{equation}\label{ch08:kp-condition}
 \sum_{\eta:\eta\not\sim\gamma}
       |z_\eta|e^{A_\eta}\le A_\gamma,
 \qquad \sum_\gamma |z_\gamma|e^{A_\gamma}<\infty.
\end{equation}
Here the first inequality is required for each $\gamma\in\mathcal P$.
Then the cluster series
\begin{equation}\label{ch08:log}
 \mathcal L(z)=\sum_{n\ge1}\frac1{n!}
       \sum_{\gamma_1,\ldots,\gamma_n}
        \varphi^T(\boldsymbol\gamma)\prod_{i=1}^n z_{\gamma_i}
\end{equation}
converges absolutely and $Z_{\rm pol}(z)=e^{\mathcal L(z)}$.
For a fixed label $\gamma_1$, its absolute rooted series satisfies
\begin{equation}\label{ch08:root}
 \sum_{n\ge1}\frac1{(n-1)!}
 \sum_{\gamma_2,\ldots,\gamma_n}
 |\varphi^T(\boldsymbol\gamma)|\prod_{i=2}^n|z_{\gamma_i}|
 \le e^{A_{\gamma_1}}.
\end{equation}
\end{theorem}

\begin{proof}
This is the hard-core specialization of the cited convergence theorem:
the Mayer function is $-1$ on incompatible pairs and zero otherwise.
There is no need to assume $Z_{\rm pol}>0$ before taking the logarithm.
The series defines the branch continued from zero activity. For real
activities it is real, and the exponential identity then gives
$Z_{\rm pol}>0$. For complex activities the conclusion is nonvanishing.
\end{proof}

The connected graph sum must be taken before its absolute value. Bounding
the sum of the absolute values of all connected graphs would discard the
cancellation on which convergence depends. The standard tree-graph
inequality is one way to establish \eqref{ch08:root}; we use the result
rather than repeat its combinatorial proof.

\section{A spatially local bound with a geometric tail}
\label{ch08:tail}

The site norm from \eqref{ch08:site-norm} gives a particularly transparent
criterion. Suppose that, for some $\alpha\ge0$, $\delta>0$ and
$\gamma_0>0$,
\begin{equation}\label{ch08:reserve}
 \|z\|_{\alpha+\delta}\le\gamma_0,
 \qquad 0<q<1,\qquad h\gamma_0\le q\delta.
\end{equation}
The additional decay $e^{-\delta c(X)}$ controls the sum over polymers
whose supports meet $X$. The parameter $q$ measures
smallness and controls the number of occurrences in a cluster.

Write $X(\boldsymbol\gamma)=\bigcup_i X_{\gamma_i}$. A direct
weighted application of Theorem~\ref{ch08:kp} yields, for every site $x$
and integer $k\ge1$,
\begin{equation}\label{ch08:weighted-tail}
 \sum_{n\ge k}\frac1{n!}
 \sum_{\boldsymbol\gamma:x\in X(\boldsymbol\gamma)}
 |\varphi^T(\boldsymbol\gamma)|
 \prod_{i=1}^n|z_{\gamma_i}|
 e^{\alpha c(X(\boldsymbol\gamma))}
 \le\gamma_0 q^{k-1}.
\end{equation}
This form is slightly stronger than the bound with an additional
$(1-q)^{-1}$ obtained by separately summing estimates at each occurrence
number. Either form is sufficient for the applications.

To prove \eqref{ch08:weighted-tail}, apply the rooted bound
\eqref{ch08:root} to the rescaled activities
$\widetilde z_\gamma=q^{-1}|z_\gamma|e^{\alpha c(X_\gamma)}$
and take $A_\gamma=\delta c(X_\gamma)$. By
\eqref{ch08:overlap} and \eqref{ch08:cover-properties}, the left side of
\eqref{ch08:kp-condition} is bounded by
$q^{-1}|X_\gamma|\gamma_0\le\delta c(X_\gamma)$.
For a nonzero connected coefficient the union weight is at most the
product of the support weights. Designating one occurrence whose support
contains $x$ changes $1/n!$ to $1/(n-1)!$ and gives an upper bound;
the rooted sum is at most $\gamma_0/q$. Restoring the factor $q^n$
on terms with $n\ge k$ gives \eqref{ch08:weighted-tail}. Nonempty supports
and finite $\Lambda$ verify the total summability required by the theorem.

More generally, let $\mathcal Q\subset\mathcal P$ be a distinguished
set of labels with
\[
 \sum_{\gamma\in\mathcal Q}|z_\gamma|
                  e^{(\alpha+\delta)c(X_\gamma)}\le b_{\mathcal Q}.
\]
Replacing the condition $x\in X(\boldsymbol\gamma)$ by
``at least one occurrence belongs to $\mathcal Q$'' in
\eqref{ch08:weighted-tail} gives $b_{\mathcal Q}q^{k-1}$ on the right.
This is useful when a cluster has to meet a prescribed insertion.
An insertion selected in this way is a distinguished occurrence in the
combinatorial expansion, not yet a renormalized quantum field.

For a family satisfying \eqref{ch08:reserve} with the same constants in
every volume, the right-hand side of \eqref{ch08:weighted-tail} is uniform
in $|\Lambda|$. Summing the $k=1$ estimate over $x$ bounds $|\mathcal L(z)|$ by
$|\Lambda|\gamma_0$. Thus extensive free-energy bounds coexist with
volume-independent estimates on connected terms meeting a fixed site.
It is the latter estimates that preserve locality in a renormalization
step.

\section{Deleting polymers and controlling local changes}
\label{ch08:deletion}

Recall that $Z_{\rm pol}$ is the gas partition function
\eqref{ch08:gas}, and $X_\gamma$ is the full support of label $\gamma$.
For a set $Y\subset\Lambda$, delete all labels whose supports meet $Y$:
put $z^Y_\gamma=z_\gamma$ if $X_\gamma\cap Y=\varnothing$ and zero
otherwise. Under \eqref{ch08:reserve}, both cluster series converge and
\begin{equation}\label{ch08:deletion-identity}
 \frac{Z_{\rm pol}(z^Y)}{Z_{\rm pol}(z)}
  =e^{-\mathcal L_Y(z)},\qquad
 \mathcal L_Y(z)=\sum_{n\ge1}\frac1{n!}
 \sum_{\substack{\boldsymbol\gamma\in\mathcal P^n\\
 X(\boldsymbol\gamma)\cap Y\ne\varnothing}}
 \varphi^T(\boldsymbol\gamma)\prod_{i=1}^n z_{\gamma_i}.
\end{equation}
This is simply subtraction of two absolutely convergent logarithms.
Its local bound is
\begin{equation}\label{ch08:deletion-bound}
 |\mathcal L_Y(z)|\le |Y|\gamma_0,
\end{equation}
and clusters with at least $k$ occurrences contribute at most
$|Y|\gamma_0q^{k-1}$. Counting a cluster more than once when it meets
several sites of $Y$ only increases the upper bound.

The same mechanism controls changes of activities in a bounded region.
If two families of activities agree away from $Y$, their logarithmic
cluster terms cancel unless a cluster contains a changed activity.
A quantitative Lipschitz estimate additionally needs a bound on the
difference of those activities, or a uniformly admissible interpolation
between them. Agreement of supports alone gives locality of the
difference, not its smallness.

\section{Parameters, derivatives and connected insertions}
\label{ch08:parameters}

Suppose $z_\gamma(u)$ is holomorphic in a complex parameter $u$ on a
disc and that \eqref{ch08:reserve} holds uniformly on each of its compact
subdiscs. The cluster series then converges normally there. Consequently
its sum is holomorphic and can be differentiated term by term. If the
bound for a local cluster sum is $M$ on $|u-u_0|<\rho$, Cauchy's estimate
gives
\begin{equation}\label{ch08:cauchy}
 |\partial_u^r\mathcal L_Y(u_0)|\le r!\rho^{-r}M.
\end{equation}
There is no separate combinatorial proof for each derivative order.
Several parameters are treated on a polydisc, with the product of the
corresponding Cauchy factors. Only derivatives of fixed finite order are
asserted unless additional uniform estimates are supplied.

One can introduce formal sources for bounded quantities $F_j$ by the
finite-volume generating function
\[
 Z(\boldsymbol t)=Z(0)\,
       \mathbb E\exp\!\left(\sum_{j=1}^r t_jF_j\right).
\]
The mixed derivative of $\log Z$ at zero is the connected correlation
(or cumulant) of the $F_j$. For example,
\[
 \left.\partial_{t_1}\partial_{t_2}\log Z(\boldsymbol t)
 \right|_{\boldsymbol t=0}
 =\mathbb E(F_1F_2)-\mathbb E(F_1)\mathbb E(F_2).
\]
The cluster representation expresses such derivatives through clusters
carrying the distinguished source occurrences; see
\cite[\S3]{Ueltschi2004}. To use this representation in the gauge theory,
the source-deformed activities must satisfy a uniform criterion on an
actual complex neighbourhood of zero. Boundedness of the undeformed gas
does not imply that estimate.

There is a second way to insert observables: differentiate the unnormalized
density and expand the result directly. This is preferable if some activity
can vanish, because division by that activity is then meaningless.
Bounds are assigned to the inserted activity itself. The usual mnemonic
``differentiate the logarithm of each activity'' is not a valid definition
in that situation.

\section{Large-field constraints and exact reassembly}
\label{ch08:forced}

An effective gauge density includes regions where a small-field expansion
is not available. Fix a configuration of the coarse variables for the
moment, and let $D\subset\Lambda$ be the set of sites declared to be in
the large-field region. In addition to the regular activities
$z_\gamma$, introduce large-field activities $K_\eta$ with supports
$Y_\eta$ and nonempty sets $I_\eta\subset D\cap Y_\eta$.
The set $I_\eta$ records the large-field sites assigned to that activity.
Let $\mathfrak C(D)$ consist of the allowed finite families $B$ of such
labels with disjoint supports and
\[
 \bigcup_{\eta\in B} I_\eta=D.
\]
The exact rules for eligibility are part of the definition of an effective
density; they may exclude some otherwise possible covers. When $D$ is
empty, the empty family is permitted. Figure~\ref{ch08:fig-exact-cover}
shows two covers of the same set \(D\) and how the exclusion rule can fail
even at a site outside \(D\).

% !TEX root = ../monograph.tex
\begin{figure}[htbp]
\centering
\begin{tikzpicture}[x=1cm,y=.70cm,font=\small,
  support/.style={draw=ymInk,line width=.85pt,rounded corners=2pt},
  lattice/.style={draw=ymGrid,thin}]
  \foreach \offset/\case/\titletext in
    {0/1/One support,4.4/2/Two disjoint supports,8.8/3/Excluded pair}{
    \begin{scope}[xshift=\offset cm]
      \ifcase\case\or
        \draw[support,fill=ymInk!7] (.67,.65) rectangle (2.33,2.35);
      \or
        \draw[support,fill=ymInk!7] (.67,.65) rectangle (1.33,2.35);
        \draw[support,fill=ymInk!7] (1.67,.65) rectangle (2.33,2.35);
      \or
        \draw[support,fill=ymInk!7] (.67,.65) -- (1.33,.65) --
          (1.33,2.68) -- (2.36,2.68) -- (2.36,3.36) --
          (.67,3.36) -- cycle;
        \draw[support] (1.73,.65) rectangle (2.27,3.25);
      \fi
      \foreach \x in {0,...,3}
        \draw[lattice] (\x,.44) -- (\x,3.48);
      \foreach \y in {1,2,3}
        \draw[lattice] (-.15,\y) -- (3.15,\y);
      \draw[ymAccent,line width=1.35pt] (1,1) rectangle (2,2);
      \foreach \x in {1,2}
        \foreach \y in {1,2}
          \fill (\x,\y) circle[radius=2.2pt];
      \ifnum\case=3
        \draw[ymInk,line width=.8pt,fill=white] (2,3) circle[radius=2.5pt];
      \fi
      \node at (1.5,-.04) {\titletext};
    \end{scope}
  }
\end{tikzpicture}
\caption[Exact covers of a large-field set]{Exact covers in a coordinate plane. One plaquette exceeds the
large-field threshold. Its four vertices (dots) form \(D\).
The outlined supports in either of the first two panels cover \(D\)
exactly, assuming the additional eligibility tests hold. In the middle
panel the indicators in both supports read the four links of the same
highlighted plaquette. In the last panel the prescribed large-field sets
still partition \(D\), but the supports share the circled site outside \(D\).}
\label{ch08:fig-exact-cover}
\end{figure}


For the simple overlap incompatibility rule, the constrained partition
sum has the form
\begin{equation}\label{ch08:constrained}
 Z_D=\sum_{B\in\mathfrak C(D)}
       \left(\prod_{\eta\in B}K_\eta\right)
       Z_{\rm pol}(z^{Y(B)}),\qquad
 Y(B)=\bigcup_{\eta\in B}Y_\eta.
\end{equation}
Applying \eqref{ch08:deletion-identity} gives the exact factorization
\begin{equation}\label{ch08:reassembly}
 Z_D=e^{\mathcal L(z)}
       \sum_{B\in\mathfrak C(D)}
         \left(\prod_{\eta\in B}K_\eta\right)
         e^{-\mathcal L_{Y(B)}(z)}.
\end{equation}
For finite sums this is algebra; the same identity holds for absolutely
convergent sums. It is the basic reason that removing regular polymers
near a large-field region can be absorbed into modified activities.

The expression $\mathcal L_{Y(B)}$ need not split as a sum over individual
members of $B$. A connected regular cluster can meet two large-field
supports. Such a cluster couples them, and the modified activity belongs
on their union together with the cluster support. Grouping by connected
components is the appropriate reorganization. An arbitrary product of
independently corrected large-field activities would omit these terms.

Formula \eqref{ch08:reassembly} does not by itself bound the sum over
$B$. That estimate requires suppression of large-field activities by the
action, bounds on the number and geometry of eligible supports, and enough
additional exponential support decay to absorb the regular clusters attached to them.
These are analytic inputs to the nonlinear renormalization step. The
abstract theorem bounds the regular cluster sum once those inputs have
been established; it does not establish them.

Two support notions must be distinguished in this application. The set
used in an exclusion or covering rule need not contain every link read
by an activity. For example, a large-field condition at a boundary site can
inspect a plaquette extending outside its assigned covering set. Disjoint
covering sets therefore need not imply disjoint dependence on links.
The abstract gas above needs the overlap implication
\eqref{ch08:overlap} for its chosen exclusion supports; it does not assume
statistical independence. When a local integration or factorization
argument also requires full dependence supports, those enlarged sets
and their geometric estimates must be specified separately. One must not
silently replace the covering sets by enlarged sets in an exact covering
identity.

\section{What an application must supply}
\label{ch08:application}

The reusable conclusions of this chapter are the connected expansion
\eqref{ch08:log}, the local tail \eqref{ch08:weighted-tail}, parameter
differentiation under a uniform complex-domain bound, and exact removal
and reassembly formulas. They are statements about a specified family
of activities. To apply them after a block integration one must give:
\begin{enumerate}
 \item an exact expansion of the integrated density, including the
 scalar factor and every large-field sector;
 \item supports containing every link and cutoff condition on which an
 activity depends, and a definite incompatibility relation;
 \item a site-norm bound with a strictly positive support-weight reserve,
 satisfying \eqref{ch08:reserve} uniformly in the period;
 \item corresponding estimates for the source-deformed or inserted
 activities whenever a connected observable is required;
 \item an absolutely summable bound for the constrained large-field
 families and the clusters attached to them.
\end{enumerate}
These requirements explain why the cluster expansion is not the main
obstruction in the construction. The general summation mechanism is
already available in the literature. What must be proved for the Wilson
model is that exact blocking, stationary integration, the large-field
decomposition and the chosen observable normalization produce activities
obeying those requirements on the same trajectory.
